\section*{Erd\H{o}s Problem 1080}

\paragraph{Statement and result.}
There are bipartite graphs of arbitrarily large order \(n\), with one part of size exactly \(\lfloor n^{2/3}\rfloor\), which have no \(6\)-cycle and have \(e(G)/n\to\infty\). In fact the construction below gives \(e(G)\gg n^{16/15}\). This disproves the proposed linear edge bound.

\paragraph{A finite Hermitian incidence geometry.}
Let \(q\) tend to infinity through odd primes and let \(F=\mathbb F_{q^2}\). On \(F^4\) use the nondegenerate Hermitian form
\[
 h(x,y)=\sum_{i=1}^4 x_i y_i^q.
\]
The points \(\mathcal P\) are the one-dimensional \(F\)-subspaces spanned by nonzero vectors \(x\) with \(h(x,x)=0\). The lines \(\mathcal L\) are the two-dimensional totally isotropic subspaces: \(h\) vanishes on every pair of their vectors. Join a point to a line when it is contained in that line.
We use the standard finite-field facts that \(F\) exists, its multiplicative group is cyclic, and the norm \(z\mapsto z^{q+1}\) maps onto \(\mathbb F_q^\times\), with exactly \(q+1\) preimages of each nonzero value.

\paragraph{Counting points and lines.}
Let \(N_d\) count vectors in \(F^d\) whose coordinate norms sum to zero, including the zero vector. Considering the first \(d-1\) coordinates, the final norm equation has one solution if its right side is zero and \(q+1\) otherwise. Hence
\[
 N_1=1,\qquad
 N_d=(q+1)q^{2d-2}-qN_{d-1}.
\]
In particular
\[
 N_2=q^3+q^2-q,\qquad
 N_4=q^7+(q-1)q^3.
\]
Every point has \(q^2-1\) nonzero representatives, so
\[
 v:=|\mathcal P|=\frac{N_4-1}{q^2-1}
 =(q^3+1)(q^2+1).
 \tag{1}
\]
Each line contains \(q^2+1\) points. If \(P\) is a point, the quotient \(P^\perp/P\) is a nondegenerate Hermitian space of dimension two. Its isotropic one-spaces correspond exactly to the lines through \(P\), and their number is
\[
 \frac{N_2-1}{q^2-1}=q+1.
 \tag{2}
\]
For completeness, a nondegenerate Hermitian form over \(F\) has an orthogonal basis with all diagonal entries equal to one: first choose a vector of nonzero norm, rescale by norm-surjectivity, and continue on its orthogonal complement. Such a vector exists because if every \(h(x,x)\) vanished, substituting \(x+y\) and \(x+\lambda y\) for \(\lambda\notin\mathbb F_q\) would force every \(h(x,y)\) to vanish. This justifies applying the \(N_2\) count to the quotient. Double-counting incidences now gives
\[
 |\mathcal L|=\frac{v(q+1)}{q^2+1}
 =(q^3+1)(q+1).
 \tag{3}
\]

\paragraph{Why a \(6\)-cycle is impossible.}
Suppose distinct points \(P_1,P_2,P_3\) and distinct lines formed a \(6\)-cycle. Every pair of these points lies on one of its lines, so their vector representatives are pairwise orthogonal and individually isotropic. If the three representatives were linearly independent, their span \(U\) would be a three-dimensional totally isotropic subspace. But \(U\subseteq U^\perp\) in a nondegenerate four-dimensional space implies
\(2\dim U\leq4\), a contradiction. If the representatives are dependent, they span a two-dimensional subspace, and all three lines must equal that span: two distinct one-spaces determine their two-dimensional span uniquely. This again contradicts the cycle. Thus the incidence graph is \(C_6\)-free.

\paragraph{Choosing the smaller part and the exact order.}
Set
\[
 M=\left\lceil2v^{2/3}\right\rceil,\qquad
 n=\left\lceil M^{3/2}\right\rceil.
\]
As \(q\to\infty\), we have \(v\asymp q^5\), \(M\asymp q^{10/3}\), and \(|\mathcal L|\asymp q^4\), so \(M\leq|\mathcal L|\) for all sufficiently large \(q\). Select any \(M\) lines and retain all \(v\) points and their incidences with those lines. Since
\[
 \frac{n}{v}\longrightarrow2^{3/2},\qquad
 \frac{M}{v}\longrightarrow0,
\]
we also have \(n-M\geq v\) eventually. Add \(n-M-v\) isolated vertices to the point part. The resulting graph has total order \(n\).
For every positive integer \(M\), the interval
\([M^{3/2},(M+1)^{3/2})\) has length greater than one. Therefore
\[
 M^{3/2}\leq n<(M+1)^{3/2},
 \qquad \lfloor n^{2/3}\rfloor=M.
\]
The selected-line part has exactly the required size. Each selected line still has all its \(q^2+1\) neighbours, so
\[
 e(G)=M(q^2+1)\asymp q^{16/3}
 \asymp n^{16/15},\qquad
 \frac{e(G)}n\asymp q^{1/3}\longrightarrow\infty.
\]
Deleting lines and adding isolated vertices cannot introduce a cycle. This completes the disproof.

\paragraph{Attribution and assessment.}
The geometry used is the classical Hermitian generalized quadrangle, and the argument explains the finite-field counts and the absence of \(6\)-cycles directly. The exponent \(16/15\) is recorded for this problem in the work of Lazebnik--Ustimenko--Woldar; no novelty is claimed for this reconstruction. Apart from the explicitly stated standard finite-field facts, the counting and graph argument are proved above. Source statement and bibliographic record checked 24 September 2026: \url{https://www.erdosproblems.com/1080}. No local formal verification is claimed.

\paragraph{Completion estimate.}
\textbf{COMPLETION: 100\%}
